% =============================================================================
\section{Justification du choix final}
\label{sec:justification}
% =============================================================================

Le choix de CMA-ES comme méthode finale repose sur trois arguments convergents :

\begin{enumerate}
  \item \textbf{Supériorité théorique pour le problème considéré.}
    CMA-ES est reconnu comme référence
    (\emph{state-of-the-art}) pour l'optimisation continue en boîte noire
    jusqu'à quelques centaines de dimensions~\cite{hansen2016}.  Notre problème
    ($d = 8$--$32$, pas de gradient, multimodal) tombe exactement dans ce régime.

  \item \textbf{Adaptation aux corrélations entre variables.}
    Les coordonnées $(x_i, y_i)$ des points de contrôle successifs sont
    naturellement corrélées : un bon chemin implique un alignement progressif
    des points, ce que $\mathbf{C}$ capture automatiquement.  GA (SBX) et DE
    (croisement binomial) travaillent composante par composante et ne
    bénéficient pas de cette adaptation.

  \item \textbf{Résultats expérimentaux.}
    CMA-ES domine ou égale DE et GA sur les 4 instances, avec des écarts
    particulièrement marqués sur \texttt{prob1} ($d = 32$) et \texttt{prob3}
    (slalom corrélé).
\end{enumerate}
